\documentclass[11pt]{article}
\usepackage[margin=1in]{geometry}
\usepackage[T1]{fontenc}
\usepackage[utf8]{inputenc}
\usepackage{amsmath,amssymb,amsthm}
\usepackage{enumitem}

\title{ICPC World Finals 2022\\Y. Compression}
\author{}
\date{}

\begin{document}
\maketitle

\section*{Problem Summary}

We are given a binary string and may repeatedly apply the allowed compression operation. We must
output a shortest string that can be reached.

\section*{Key Observations}

\begin{itemize}[leftmargin=*]
    \item The first character never changes.
    \item The last character never changes.
    \item If the original string contains both \texttt{0} and \texttt{1}, then no sequence of operations can
    delete all occurrences of one of the two characters.
    \item These three invariants already determine the unique shortest answer.
\end{itemize}

\section*{Algorithm}

\begin{enumerate}[leftmargin=*]
    \item If all characters are equal, output that single character.
    \item Otherwise, both bits occur in the string.
    \item If the first and last characters are different, output exactly those two characters.
    \item If the first and last characters are equal, output
    \[
    \text{first} \;+\; \text{opposite bit} \;+\; \text{last}.
    \]
\end{enumerate}

\section*{Correctness Proof}

We prove that the algorithm returns the correct answer.

\paragraph{Lemma 1.}
Every reachable string has the same first character and the same last character as the original
string.

\paragraph{Proof.}
The allowed compression never changes the two endpoints of the current string, so by induction
over all operations the first and last characters remain invariant. \qed

\paragraph{Lemma 2.}
If the original string contains both \texttt{0} and \texttt{1}, then every reachable string also contains
both \texttt{0} and \texttt{1}.

\paragraph{Proof.}
This is the third key invariant of the operation: it is impossible to erase all occurrences of one
character while the other remains. Therefore any reachable string from a mixed binary string must
still contain both bits. \qed

\paragraph{Lemma 3.}
The string produced by the algorithm is reachable and no shorter reachable string exists.

\paragraph{Proof.}
If the string is constant, a single-character answer is clearly optimal.

Now assume both bits occur. By repeatedly compressing inside equal runs, we can first transform
the string into an alternating one. After that, repeated compressions of the front remove two
characters at a time while preserving the same endpoint characters. Thus we can always reduce to
the shortest alternating string with the same endpoints.

If the endpoints differ, the shortest such alternating string is exactly the length-$2$ string formed
by those endpoints. If the endpoints are equal, a length-$2$ string is impossible, and because both
bits must remain present by Lemma 2, the shortest possibility is the length-$3$ string
\texttt{first opposite first}. Hence the algorithm's answer is reachable and optimal. \qed

\paragraph{Theorem.}
The algorithm outputs a shortest reachable string.

\paragraph{Proof.}
By Lemma 1 and Lemma 2, any reachable optimum must satisfy the same endpoints and, when the
input is mixed, must contain both bits. Lemma 3 shows that the algorithm constructs exactly the
shortest string satisfying those necessary conditions. Therefore the answer is correct. \qed

\section*{Complexity Analysis}

We only scan the string once to test whether all characters are equal. The running time is $O(n)$
and the memory usage is $O(1)$.

\section*{Implementation Notes}

\begin{itemize}[leftmargin=*]
    \item Once the three cases above are identified, the answer can be printed directly without
    simulating any compression steps.
\end{itemize}

\end{document}
